\documentclass[preprint,12pt]{elsarticle}
\usepackage[T1]{fontenc}
\usepackage[utf8]{inputenc}
\usepackage[bahasa]{babel}
\usepackage{graphicx}
\usepackage{booktabs}
\usepackage{tabularx}
\usepackage{array}
\usepackage{amsmath}
\usepackage{longtable}
\usepackage[hidelinks]{hyperref}
\usepackage{float}
\newcolumntype{L}[1]{>{\raggedright\arraybackslash}p{#1}}
\newcolumntype{Y}{>{\raggedright\arraybackslash}X}
\newcommand{\E}[1]{E-#1}
\newcommand{\G}[1]{G#1}
\newcommand{\mAP}{mAP}

\journal{Computers and Electronics in Agriculture}
\title{Apakah Kanal Keempat Depth Meningkatkan Deteksi Tandan Buah Segar
Kelapa Sawit oleh YOLO? Hasil 32 Eksperimen}
\author{Muhammad Zainal Muttaqin dan Fatma Indriani}
\date{4 Agustus 2026}

\begin{document}
\begin{frontmatter}
\maketitle

\begin{abstract}
Laporan ini menjawab satu pertanyaan. Detektor YOLO membaca tiga kanal RGB.
Kami menambahkan depth sebagai kanal keempat. Kami mengukur apakah deteksi
meningkat, tetap sama, atau justru menurun.

Jawabannya, depth tidak meningkatkan deteksi. Fusi awal merugikan pada dua
dari tiga seed. Fusi menengah dan fusi akhir menggeser skor lebih kecil
daripada derau eksperimen itu sendiri. Seluruh dua belas selang kepercayaan
berpasangan memuat nol. Kanal kendali berisi derau acak menggeser skor pada
orde besaran yang sama dengan kanal depth sungguhan.

Kami juga mengukur apakah implementasi depth itu sendiri sudah benar. Pada
mulanya tidak benar. Metadata sensor mengklaim penyelarasan ke kamera warna,
tetapi buffer tetap berada di grid kamera depth. Kami memperbaiki galat ini
dengan reproyeksi penuh. Hasil negatif di atas berasal dari pipeline yang sudah
diperbaiki, bukan dari pipeline yang cacat.

Tiga puluh dua eksperimen bernomor (\E{001} sampai \E{032}) dan sembilan
gerbang (\G{0} sampai \G{8}) memikul bukti. Laporan ini mendaftar setiap
eksperimen beserta hasilnya dalam satu baris.
\end{abstract}

\begin{keyword}
deteksi tandan buah segar \sep kelapa sawit \sep RGB-D \sep masukan empat kanal
\sep Orbbec \sep YOLO26 \sep ablasi fusi
\end{keyword}
\end{frontmatter}

\section{Pertanyaannya}
\label{sec:question}

Detektor YOLO membaca citra dengan tiga kanal, yaitu merah, hijau, dan biru.
Kamera depth dapat memasok kanal keempat. Kanal keempat itu memuat jarak dari
kamera ke setiap piksel.

Gagasannya sederhana. Dua tandan buah dapat berwarna sama dan dapat bertindih
di dalam citra. Dua tandan yang sama itu dapat berdiri pada jarak berbeda dari
kamera. Depth dapat memisahkan keduanya ketika warna tidak dapat.

Laporan ini mengukur satu hal:

\begin{itemize}
\item Apakah kanal keempat menaikkan akurasi deteksi?
\item Apakah kanal keempat tidak mengubah apa pun?
\item Apakah kanal keempat menurunkan akurasi karena menambah derau?
\end{itemize}

Satu-satunya cara menjawab pertanyaan ini adalah menjalankan eksperimennya dan
membaca angkanya. Laporan ini memuat angka-angka tersebut.

\section{Jawabannya}
\label{sec:answer}

\subsection{Jawaban singkat}

Kanal keempat tidak menaikkan akurasi deteksi pada regime yang diuji. Dua dari
tiga seed fusi awal memburuk. Fusi menengah dan fusi akhir menggeser skor lebih
kecil daripada ragam antar-seed pada model yang sama.

\begin{table}[H]
\centering
\caption{Hasil matriks fusi (\E{032}). Setiap nilai adalah perubahan
\mAP50 test terhadap model RGB-saja dengan seed yang sama.}
\label{tab:answer}
\begin{tabularx}{\columnwidth}{L{0.20\columnwidth}L{0.34\columnwidth}L{0.14\columnwidth}Y}
\toprule
\textbf{Titik fusi} & \textbf{Perubahan per seed} & \textbf{Rerata} & \textbf{Hasil}\\
\midrule
Awal & $-0.0120$, $+0.0234$, $-0.0017$ & $+0.0032$ & Tanda berubah-ubah. Tidak ada kenaikan.\\
Menengah & $+0.0096$, $+0.0212$, $+0.0110$ & $+0.0139$ & Positif pada 3 dari 3 seed. Selangnya memuat nol.\\
Akhir & $-0.0056$, $+0.0070$, $+0.0102$ & $+0.0039$ & Tanda berubah-ubah. Tidak ada kenaikan.\\
Derau acak & $-0.0130$, $+0.0025$, $-0.0081$ & $-0.0062$ & Lengan kendali. Tidak memuat geometri.\\
\bottomrule
\end{tabularx}
\end{table}

Seluruh dua belas selang kepercayaan berpasangan memuat nol. Fusi menengah
memiliki rerata terbesar dan bertanda positif pada ketiga seed. Ini indikasi.
Ini bukan peningkatan yang terkonfirmasi.

\subsection{Mengapa lengan derau penting}

Lengan derau adalah jawaban paling langsung atas pertanyaan ini. Kami mengganti
kanal depth dengan bilangan acak. Lengan derau kemudian menggeser skor pada
orde besaran yang sama dengan ketiga lengan depth. Tidak ada satu pun kontras
yang mengecualikan nol.

Pengukuran ini tidak membuktikan bahwa depth dan derau setara. Kami tidak
menjalankan uji ekuivalensi. Pengukuran ini menunjukkan bahwa nilai depth tidak
memperlihatkan informasi yang dapat dimanfaatkan pada regime ini, karena kanal
tanpa geometri menghasilkan perubahan berukuran sama.

\subsection{Apa yang tidak dikatakan jawaban ini}

Jawaban ini berlaku untuk regime yang diuji. Regime itu adalah YOLO26n, 640
piksel, 150 epoch, satu split pohon, tiga seed, dan pelatihan dari awal. Kami
tidak membuktikan bahwa depth dan RGB setara. Kami tidak mengukur model fusi
menengah yang terlatih awal, karena tidak ada checkpoint COCO yang kompatibel
untuk topologi dua cabang. Kami tidak mengukur perkebunan yang lebih luas atau
model kamera kedua.

\section{Apakah implementasi depth sudah benar?}
\label{sec:implementation}

Hasil negatif hanya berguna apabila implementasinya benar. Kami mengukur
implementasinya lebih dahulu. Di dalamnya terdapat lima galat. Kami memperbaiki
kelimanya sebelum membaca angka final. Empat galat pertama adalah milik pipeline
depth dan berasal dari audit \E{022}. Galat kelima adalah asimetri evaluator,
yang diukur \E{025} sebagai gerbang tersendiri.

\begin{table}[H]
\centering
\caption{Galat yang ditemukan pada pipeline depth beserta perbaikannya
(\E{022} dan \E{025}).}
\label{tab:defects}
\begin{tabularx}{\columnwidth}{L{0.26\columnwidth}Y}
\toprule
\textbf{Galat} & \textbf{Perbaikan dan efek terukur}\\
\midrule
Registrasi keliru & Sidecar menyatakan \texttt{alignedTo: color}, tetapi buffer
tetap berada di grid kamera depth. Resize langsung meleset 29,3 piksel pada
median dan 61 piksel pada maksimum. Reproyeksi penuh menaikkan diagnostik
mutual information sebesar $+0.0306$ bit, dengan selang 95 persen
$[0.0260, 0.0354]$. Mutual information adalah diagnostik registrasi. Ia bukan
bukti penyelarasan geometris, dan bukan bukti manfaat di hilir.\\
Kalibrasi dipakai bersama & Dua unit kamera memiliki intrinsik berbeda.
Pipeline kini membaca kalibrasi dari sidecar setiap berkas.\\
Rentang depth keliru & Rentang bawaan 0,3 sampai 8,0 m tidak cocok untuk sensor
ini. Audit piksel \E{022} tidak menemukan piksel di bawah 0,3 m, dan menemukan
10,07 persen piksel di atas 8 m. Rentang terkoreksi dari train saja adalah 0,8
sampai 15,0 m.\\
Donor lintas-split dan state acak bersama & Hasil menguntungkan yang pertama
memakai citra donor dari split lain dan state bilangan acak bersama. Hasil itu
kami tarik.\\
Asimetri evaluator & Kedua lengan menghasilkan kerapatan deteksi berbeda 2,44
kali. Ini menimbulkan offset 0,0156, cukup besar untuk membalik tanda selisih
depth dikurangi RGB. Seluruh pekerjaan berikutnya memakai satu evaluasi
berpasangan \texttt{pycocotools}.\\
\bottomrule
\end{tabularx}
\end{table}

Hasil depth yang pertama tampak positif. Hasil itu berasal dari pipeline yang
cacat. Selang kepercayaannya sendiri sudah memuat nol, yaitu
$[-0.0215,\allowbreak +0.0632]$. Hasil itu kami tarik.
Setiap angka pada Bagian~\ref{sec:answer} berasal dari pipeline yang sudah
diperbaiki.

\subsection{Tiga jenis depth, dan alasan ketiganya diukur}

\begin{table}[H]
\centering
\caption{Tiga sumber depth dalam studi ini.}
\label{tab:depthkinds}
\begin{tabularx}{\columnwidth}{L{0.22\columnwidth}L{0.24\columnwidth}Y}
\toprule
\textbf{Sumber} & \textbf{Sifat} & \textbf{Alasan diukur}\\
\midrule
Pseudo-depth & Diprediksi dari citra RGB yang sama. Skala relatif. & Gratis.
Bila berhasil, kamera tidak diperlukan. Diukur pada \E{003} sampai \E{007}.\\
Depth fisik & Sensor Orbbec Y16. Milimeter metrik. & Merupakan pengukuran
independen. Diukur pada \E{022} sampai \E{032}.\\
Derau acak & Nilai acak pada posisi tensor yang sama. & Menunjukkan biaya kanal
tambahan itu sendiri. Diukur pada \E{027}, \E{029}, \E{030}, dan \E{032}.\\
\bottomrule
\end{tabularx}
\end{table}

Pseudo-depth gagal pada pengukuran langsung di \E{006}. Kontras kotak tandan
adalah 0,0089 berbanding 0,0341 untuk kendali acak berukuran sepadan. Selisih
AUC adalah 0,009. Pseudo-depth berasal dari citra RGB yang sama, sehingga
galatnya berkorelasi dengan galat RGB. Ia bukan sensor independen.

Depth fisik lolos audit berkas setelah perbaikan pada
Tabel~\ref{tab:defects}. Setelah itu ia gagal pada pengukuran manfaat.

\section{Bukti: seluruh 32 eksperimen}
\label{sec:experiments}

Setiap baris memuat satu eksperimen, pertanyaannya, dan hasilnya. Konfigurasi
lengkap, lokasi artefak, dan nama skrip tersimpan di log repositori
\texttt{experiments/EKSPERIMEN.md}.

\begin{longtable}{L{0.07\textwidth}L{0.28\textwidth}L{0.55\textwidth}}
\caption{32 eksperimen bernomor dan hasilnya.}
\label{tab:allexp}\\
\toprule
\textbf{ID} & \textbf{Pertanyaan} & \textbf{Hasil}\\
\midrule
\endfirsthead
\toprule
\textbf{ID} & \textbf{Pertanyaan} & \textbf{Hasil}\\
\midrule
\endhead
\bottomrule
\endfoot

\E{001} & Apakah statistik \texttt{class\_mismatch} mengukur ambiguitas
kematangan? & Tidak. 0 dari 7.328 tandan multi-sisi berbeda pendapat, karena
pipeline yang sama menghasilkan statistik dan tautannya sekaligus. DIPALSUKAN
sebagai alat ukur.\\

\E{002} & Dapatkah anotasi master mentah dipakai ulang secara langsung? & Tidak.
3.992 berkas mentah hanya memuat 1.352 basename unik, dan 936 di antaranya
berulang antar-folder. TIDAK KONKLUSIF, sehingga \E{015} memetakan berbasis
isi.\\

\E{003} & Dapatkah pseudo-depth menata satu video orbit? & Sebagian. 48 frame
mencakup 319,7 derajat dengan residu lingkaran 8,2 persen. Satu video tidak
cukup. PARSIAL.\\

\E{004} & Apakah itu berlaku pada enam video? & Ya untuk geometri pohon. Lima
dari enam video mencapai 270 derajat atau lebih. Satu video mencapai 41 persen.
TERKONFIRMASI hanya untuk geometri pohon.\\

\E{005} & Dapatkah pseudo-depth mengurutkan sisi pemotretan? & Ya. Seluruh 20
pohon empat sisi dan seluruh 30 pohon delapan sisi terurut benar. RMSE sudut
adalah 17,3 dan 8,5 derajat terhadap baseline acak 57,5 dan 34,4 derajat.
TERKONFIRMASI.\\

\E{006} & Apakah pseudo-depth memisahkan tandan dari latarnya? & Tidak. Kontras
0,0089 berbanding 0,0341 untuk kendali. Selisih AUC 0,009. DIPALSUKAN.\\

\E{007} & Dapatkah pseudo-depth menautkan tandan yang sama antar-sisi? & Tidak.
Koreksi statistik sederhana lebih unggul. Sapuan ambang menyentuh split test,
sehingga nilainya bersifat eksploratif. DIPALSUKAN untuk implementasi ini.\\

\E{008} & --- & Nomor cadangan. Tidak dijalankan. Slot ini tetap ditampilkan
agar sebuah rumpang tidak tampak seperti hasil nol.\\

\E{009} & Apakah B4 sulit karena kotaknya kecil? & Tidak semata-mata. 81,2
persen kotak B4 berada di pita medium COCO. Median kotak adalah 46 kali 46
piksel. PARSIAL.\\

\E{010} & Apakah B4 sulit karena tandan berdesakan? & Tidak. B4 memiliki jumlah
tetangga terendah, IoU maksimum terendah, dan kontras rendah. B2 memiliki
$\Delta E$ tinggi dan AP50 rendah, sehingga B2 adalah masalah yang berbeda.
Desakan DIPALSUKAN.\\

\E{011} & Dapatkah prapemrosesan kontras atau tekstur memulihkan B4? & Penguatan
kontras gagal. Filter Laplacian memindahkan B4 dari kelas paling tidak terpisah
menjadi kelas paling terpisah, pada AUC 0,6153. Ini mendukung kanal tekstur,
bukan otomatis kanal depth.\\

\E{012} & Apakah kekeliruan kematangan bersifat acak? & Tidak. Akurasi 52,87
persen terhadap baseline 25 persen, dan galatnya mengikuti rantai
B1--B2--B3--B4. Galat tak bersebelahan terbesar adalah 1,9 persen. Kematangan
bersifat ordinal.\\

\E{013} & Bagaimana kontrak empat kanal yang persis? & Urutan kanal
\([B,G,R,D]\), depth tidak sahih bernilai nol persis, pipeline menyimpan mask
validitas, dan dropout modalitas menjaga inferensi RGB-saja tetap mungkin.
Kontrak lengkap. Validasi perangkat masih terbuka.\\

\E{014} & Apakah detektor gagal pada lokalisasi atau pada kematangan? & \mAP50
adalah 0,7191 tanpa kematangan dan 0,5218 dengan kematangan, rasionya 72,6
persen. PROVISIONAL, karena kami belum menunjukkan JSON prediksi yang identik
per bita.\\

\E{015} & Dapatkah piksel master dipetakan dengan aman? & Ya. 3.992 dari 3.992
citra terpetakan berbasis isi. Himpunan terpetakan itu tidak melatih \E{021}.
HAMBATAN DIBUKA.\\

\E{016} & Apakah 68 persen merupakan plafon keras kematangan? & Tidak. Ketiga
jalur pengukuran berbagi satu pengklasifikasi crop dan memakai augmentasi warna
yang tidak selaras. DITARIK.\\

\E{017} & Apakah detektor dua tahap membantu? & Tidak. Ia lebih buruk daripada
baseline satu tahap pada setiap kelas dan pada kedua metrik \mAP. DIPALSUKAN
untuk implementasi ini.\\

\E{018} & Seberapa jauh lokalisasi dapat mencapai pada kondisi terbaik? &
Keterjangkauan oracle adalah 0,8834 dan 0,4702. Median IoU terbaik adalah
0,7303, dan hanya 0,0376 kotak yang melewati IoU 0,90. Ini menjelaskan mengapa
\mAP50-95 tetap sulit. Selubung deskriptif.\\

\E{019} & Apakah menambah piksel menutup jurang? & Tidak. Satu run resolusi
tinggi yang aman warna berhenti pada epoch 41 dan tidak mencapai sasaran. Kami
mencoret resolusi semata sebagai langkah berikutnya.\\

\E{020} & Apakah transformer tanpa NMS membantu? & RT-DETR-L memperbaiki keempat
kelas, dengan kenaikan terbesar pada B4. Banyak faktor berubah bersamaan,
sehingga ini hasil daftar-pendek.\\

\E{021} & Detektor RGB mana yang menjadi kendali yang adil? & RF-DETR-L. \mAP50
test adalah 0,6038 dan \mAP50-95 test adalah 0,2770. Selisih berpasangan
terhadap RT-DETR-L adalah $+0.0255$, dengan selang 95 persen
$[0.0104, 0.0408]$. Tolok ukur RGB terbaik yang teramati.\\

\E{022} & Apakah berkas depth fisik mengukur apa yang diklaim metadatanya? &
Tidak. Lihat Tabel~\ref{tab:defects}. Galat registrasi, rentang, donor, dan
state acak sudah diperbaiki. Delta positif pertama DITARIK.\\

\E{023} & --- & Alias eksekusi. Matriks yang dirancang ulang berjalan satu kali,
sebagai \E{032}. Ini bukan pengukuran kedua.\\

\E{024} & Dapatkah identitas lintas-sisi diukur tanpa proksi sirkular? & Ya,
dari graf anotasi. Tumbukan terpusat pada B1--B2 dan B2--B3, tanpa tumbukan
B1--B3. B4 memiliki nol deteksi dua sisi, sehingga dayanya terbatas.\\

\E{025} & Apakah metrik trainer dan \texttt{pycocotools} sepakat? & Tidak. Kedua
lengan berbeda kerapatan deteksi 2,44 kali, yang menciptakan offset 0,0156.
Offset itu cukup besar untuk membalik tanda. Satu evaluator kini mengikat
seluruh pekerjaan berikutnya.\\

\E{026} & Apakah depth menstabilkan identitas lintas-sisi? & TIDAK KONKLUSIF.
Penyebutnya 82 tandan RGB dan 75 tandan RGB-D, sehingga perbandingannya tidak
sepenuhnya berpasangan. Selangnya terlalu lebar.\\

\E{027} & Apakah fusi awal membantu pada tiga seed? & Tidak. Dua dari tiga
selang bernilai negatif, dan hanya seed 42 yang positif. Kriteria manfaat
DIPALSUKAN untuk YOLO26n.\\

\E{028} & Apakah dataset yang lebih besar mengubah ukuran identitas? & Ia
menambah daya, bukan perbedaan. Lajunya 0,2329 berbanding 0,1951, selisih
$+0.0378$ dengan selang $[-0.0585, +0.1276]$. Tidak ada depth fisik pada dataset
ini.\\

\E{029} & Apakah model yang lebih besar menyelamatkan depth? & Tidak. Rerata
depth dikurangi derau adalah $+0.0124$, di bawah ambang terarsip $+0.015$.
Lengan RGB saja bervariasi 0,0759 antar-seed. Klausa penyelamatan kapasitas
DICABUT.\\

\E{030} & Apakah efeknya bergantung pada kapasitas model? & Belum jelas. Derau
dikurangi RGB adalah $+0.0032$, $+0.0184$, $-0.0325$, $-0.0533$. Tandanya
menyeberang satu kali, antara 21,9 dan 26,3 juta parameter, tetapi besarannya
tidak monoton. Hanya satu seed. Ini membuka \G{7b}.\\

\E{031} & Apakah efek split lebih besar daripada efek seed? & Ya untuk skor
absolut. Rentang \mAP absolut adalah 0,0488 antar-split berbanding 0,0321
antar-seed. Setiap skor absolut wajib menyebutkan split-nya.\\

\E{032} & Apakah titik fusi merupakan penjelasannya? & TIDAK KONKLUSIF. Lima
lengan, tiga seed, dua belas selang berpasangan, dan kedua belasnya memuat nol.
Lihat Tabel~\ref{tab:answer}.\\
\end{longtable}

\subsection{Bagaimana 32 catatan itu menjawab satu pertanyaan tersebut}

Enam catatan memberi jawaban langsung, yaitu \E{022}, \E{027}, \E{029},
\E{030}, \E{031}, dan \E{032}. Keenamnya melatih detektor dengan dan tanpa
kanal keempat lalu membandingkan skornya.

Catatan lainnya menyingkirkan penjelasan alternatif atas jawaban itu. Tanpa
catatan tersebut, pembaca dapat menolak hasil negatif ini karena alasan yang
sama sekali tidak berkaitan dengan depth:

\begin{itemize}
\item \E{001}, \E{002}, \E{015}, dan \E{024} memastikan bahwa label dan graf
identitas tidak sirkular.
\item \E{025} memastikan bahwa kedua lengan memakai satu evaluator. Sebelum
catatan ini, evaluator saja dapat membalik tanda.
\item \E{031} memastikan bahwa split disebutkan. Split menggeser skor absolut
lebih besar daripada seed.
\item \E{003} sampai \E{007} menyingkirkan pseudo-depth sebagai pengganti murah
bagi sensor.
\item \E{009} sampai \E{012} menjelaskan apa sebenarnya galat RGB itu. B4
adalah masalah kontras dan B2--B3 adalah masalah kematangan ordinal. Depth
dapat menolong yang pertama. Depth tidak dapat menolong yang kedua.
\item \E{013} membekukan kontrak kanal, sehingga run berikutnya tidak dapat
mengubah urutan kanal secara diam-diam.
\item \E{017} sampai \E{021} menyediakan kendali RGB yang kuat. Kendali yang
lemah membuat kanal keempat apa pun tampak berguna.
\end{itemize}

\section{Register gerbang \G{0} sampai \G{8}}
\label{sec:gates}

Pekerjaan depth berjalan sebagai sembilan gerbang. Gerbang adalah lubang yang
harus ditutup sebelum kesimpulan diterima. Satu gerbang lanjutan, \G{7b}, masih
terbuka.

\begin{table}[H]
\centering\small
\caption{Register gerbang. \G{0}, \G{4}, \G{6}, dan \G{7b} belum tertutup.}
\label{tab:gates}
\begin{tabularx}{\columnwidth}{L{0.13\columnwidth}L{0.24\columnwidth}L{0.14\columnwidth}Y}
\toprule
\textbf{Gerbang} & \textbf{Isi} & \textbf{Ditutup oleh} & \textbf{Status}\\
\midrule
\G{0} & Kelengkapan run dan tautan mati & \texttt{d58eae9} & \textbf{Terbuka.}
Manifes dan provenans wajib dirilis bersama sumber dan PDF.\\
\G{1} & Perbedaan metrik trainer dan \texttt{pycocotools} & \E{025} & Asimetri
terbatasi. Mekanisme lengkapnya belum tuntas.\\
\G{2} & Matriks multi-seed dalam satu protokol & \E{027}, \E{029} & Tertutup
untuk regime yang diuji. Tidak universal.\\
\G{3} & Restrukturisasi \E{022} dan penyelarasan SR-015 & \texttt{9c5d9dd},
\texttt{86f8e65} & Tertutup.\\
\G{4} & Fusi menengah & \E{032} & \textbf{Tidak konklusif.} Selangnya memuat
nol. \E{032} tidak membuktikan ekuivalensi.\\
\G{5} & Ragam antar-split & \E{031} & Diukur pada tiga split. Tidak ada hukum
populasi.\\
\G{6} & Fusi akhir & \E{032} & \textbf{Tidak konklusif.} Selangnya memuat nol.
\E{032} tidak membuktikan ekuivalensi.\\
\G{7} & Sapuan kapasitas dalam satu keluarga & \E{030} & Eksploratif. Satu
seed. Klaim kapasitas tetap terbuka.\\
\G{8} & Ukuran lintas-sisi pada dataset berdaya memadai & \E{028} & Daya RGB
bertambah. Ini bukan perbandingan depth.\\
\G{7b} & Monotonisitas kapasitas antar-seed & --- & \textbf{Terbuka.}\\
\bottomrule
\end{tabularx}
\end{table}

\subsection{Keadaan \G{7b}}

Commit \texttt{7afd274} membuka 12 run tambahan. Run tersebut mencakup yolo26m
dan yolo26l, tiga modalitas, serta seed 1337 dan 2024. Keadaan sebenarnya
adalah:

\begin{itemize}
\item 7 dari 12 run selesai berlatih, dan kurvanya terarsip.
\item Kami belum menghitung satu pun kontras berpasangan. Berkas berpasangan
hanya ada untuk seed 42.
\item Kami belum pernah memperbarui \E{030}, sehingga batasan satu seed masih
berlaku.
\end{itemize}

Karena itu, titik balik kapasitas antara 21,9 dan 26,3 juta parameter merupakan
pola satu seed. Itu bukan temuan. Lima run dan satu evaluasi berpasangan akan
menutup gerbang ini.

\section{Arti hasil ini bagi sistem lapangan}
\label{sec:consequence}

\begin{itemize}
\item Pakai jalur RGB. RF-DETR-L adalah kendali terbaik yang teramati, pada
\mAP50 test 0,6038. Kecepatannya 8,5 FPS pada GPU L4, terlalu lambat untuk
pemakaian lapangan waktu nyata tanpa optimasi.
\item Jangan menambahkan kanal depth dengan tujuan menaikkan akurasi.
Pengukuran tidak mendukung keputusan itu.
\item Apabila perangkat keras depth sudah tersedia, simpan hasil tangkapannya
sebagai cabang data yang dapat diaudit. Terapkan gerbang mutu per berkas dari
Tabel~\ref{tab:defects}.
\item Untuk galat kematangan B2--B3, gunakan loss ordinal atau kepala regresi.
\E{012} menunjuk ke arah itu. Kanal keempat tidak menyentuh batas ordinal.
\item Sasaran inventarisasi satu hitungan per tandan belum tervalidasi. Tidak
ada eksperimen dalam laporan ini yang mengukurnya.
\end{itemize}

\section{Batas jawaban ini}
\label{sec:limits}

\begin{itemize}
\item Matriks fusi hanya memakai YOLO26n, pada 640 piksel, selama 150 epoch,
pada satu split, dan dari awal.
\item Tidak ada checkpoint terlatih awal untuk topologi dua cabang, sehingga
transfer learning tidak terukur.
\item Tiga seed dan tiga split memperlihatkan ragam. Keduanya tidak menaksir
sebaran populasi.
\item Dataset depth memuat 352 pohon. Pilot RGB memuat 953 pohon.
\item Klaim ini berlaku untuk berkas Orbbec Y16 yang diaudit. Klaim ini tidak
berlaku untuk setiap kamera Orbbec atau setiap penyandian depth.
\item Kami tidak pernah melatih RF-DETR-L dengan depth fisik.
\end{itemize}

\section{Simpulan}
\label{sec:conclusion}

Kami menambahkan kanal keempat berisi depth pada detektor YOLO dan mengukur
efeknya. Pada regime yang diuji, kanal itu tidak memperbaiki deteksi. Fusi awal
merugikan pada dua dari tiga seed. Fusi menengah dan fusi akhir menghasilkan
perubahan yang juga dihasilkan kanal derau acak. Seluruh dua belas selang
berpasangan memuat nol.

Implementasi depth cacat pada awalnya, dan kami memperbaikinya. Jawaban final
berasal dari pipeline yang sudah diperbaiki, sehingga hasil negatif ini milik
sensor dan metodenya, bukan milik galat registrasi.

Jawaban ini bersyarat. Bagian~\ref{sec:limits} mendaftar syarat-syaratnya. Satu
gerbang, \G{7b}, masih dapat mengubah bagian kapasitas dari jawaban ini.
Gerbang itu memerlukan lima run dan satu evaluasi berpasangan.

\end{document}
